\section{The static base and its weak-decay end expansion}
\label{q3:static:section}

The original-data variation supplies a causal stationary field. We
verify the hypotheses of the common static-base construction, keeping
all possible ends approaching its null set. We then prove the
three-dimensional end expansion needed for conformal doubling.
The next section classifies that double and recovers the original
hypersurface through its boundary. The conformal construction follows
the method of Bunting and Masood-ul-Alam~\cite{BuntingMasoodUlAlam1987}.

\begin{proposition}[Static classification and hypersurface reconstruction]
\label{q3:static:classification}
Let the original exterior satisfy the hypotheses of
Theorem~\ref{q3:intro:rigidity}, and let \(u,X\) be the fields supplied by
Proposition~\ref{q3:variation:adjoint}.
Then the number of boundary components is $\ell=1$, and that
component is diffeomorphic to $S^2$. Moreover,
\[
N:=u^2-|X|_g^2>0
\qquad\text{on }\operatorname{int}\Mext.
\]
The metric and function
\[
h=g+N^{-1}X^\flat\otimes X^\flat,\qquad \lambda=\sqrt N
\]
form the Schwarzschild static base of mass \(m\), with its horizon
boundary attached in the regular base coordinates described below.
There is a proper smooth spacelike embedding of the original
\(\Mext\) into maximally extended Schwarzschild spacetime of mass \(m\)
whose induced metric and future second fundamental form are \(g,K\).
Its boundary is a cross-section of the future horizon if \(u|_S>0\),
and the bifurcation sphere if \(u|_S=0\).
The embedding extends as a smooth spacelike hypersurface through \(S\),
and its chosen end approaches the corresponding spatial infinity.
\end{proposition}

We record precisely the outputs of Proposition~\ref{q3:variation:adjoint}
that will be used.
Put \(\Omega=\operatorname{int}\Mext\).
The fields \(u,X\) are smooth up to the one-sided boundary,
\(u>0\) on \(\Omega\), and \(u\geq |X|_g\).
For any fixed
\[
\frac12<q_0<\min\{q,1\},
\]
their end asymptotes have an \(O_2(r^{-q_0})\) remainder and satisfy
\begin{equation}
\label{q3:static:asymptotic-field}
(u,X)\longrightarrow(b^0,b)
 =\left(\frac Em,-\frac Pm\right),
\qquad (b^0)^2-|b|^2=1,\qquad b^0>0.
\end{equation}
On \(\R\times\Omega\), the Lorentzian metric
\begin{equation}
\label{q3:static:stationary-metric}
\mathbf g=-u^2\dd t^2+
 g_{ij}(\dd x^i+X^i\dd t)(\dd x^j+X^j\dd t)
\end{equation}
has Killing field \(\xi=\partial_t\), satisfies
\(\Ric_{\mathbf g}(\xi,\cdot)=0\), and is vacuum wherever \(N>0\).
In addition,
\begin{equation}
\label{q3:static:kid-boundary}
\sym\nabla X=-uK,\qquad
X|_S=u\nu,\qquad
\partial_\nu u+K(X,\nu)=\kappa:=\frac1{4m}.
\end{equation}
For each connected component $S_i$, either $u>0$ everywhere on $S_i$
or $u=0$ everywhere on $S_i$. The alternative may initially depend on
$i$; the same positive constant $\kappa$ occurs on all components.
These statements include the null-dust alternative before vacuum has
been proved: only the contraction \(\Ric_{\mathbf g}(\xi,\cdot)=0\)
will be used at points where \(N=0\).

\subsection{The common construction under weak decay}
\label{q3:static:base-subsection}

Let \(\mathcal P=\{N>0\}\subset\Omega\). On this open set define
\begin{equation}
\label{q3:static:base-definitions}
\begin{aligned}
h&=g+N^{-1}X^\flat\otimes X^\flat,
& C&=h^{-1}=g^{-1}-u^{-2}X\otimes X,\\
\lambda&=\sqrt N,
& A&=N^{-1}X^\flat,\qquad F=\dd A .
\end{aligned}
\end{equation}
The stationary metric then has the form
\begin{equation}
\label{q3:static:threading-form}
\mathbf g=-N(\dd t-A)^2+h,\qquad
\dd V_h=\frac u{\sqrt N}\dd V_g.
\end{equation}

We apply Proposition~\ref{prop:static:common-base} to the original
closed exterior. Its completeness follows from the ambient hypothesis:
an intrinsic Cauchy sequence is Cauchy in the complete ambient manifold,
its limit remains in the closed exterior, and a smooth interior or
boundary half-ball chart gives intrinsic convergence. The exterior
has exactly one end and a compact complement of that end. Its boundary
has finitely many smooth components, all included in the application.

The end hypotheses of the common proposition require only two
differentiated bounds on the metric and stationary field. Here
\(g-\delta=O_2(r^{-q_0})\), and
Equation~\eqref{q3:static:asymptotic-field} has an
\(O_2(r^{-q_0})\) remainder, with \(2q_0>1=n-2\).
Proposition~\ref{q3:variation:adjoint} obtained that field estimate
from the original \(O_2/O_1\) data by applying the first-jet estimate
with a slightly larger input exponent. Thus no additional derivative
of the original metric or tensor is required here. The normalization
\((b^0)^2-|b|^2=1\), the positive interior lapse, causal domination,
vanishing Ricci contraction, and vacuum on \(\mathcal P\) are exactly
the outputs recorded above. Equation~\eqref{q3:static:kid-boundary}
provides the boundary law on each component with the same
\(\kappa=1/(4m)>0\). The choice between positive and zero boundary
lapse may differ from one component to another.

\begin{lemma}[The complete three-dimensional static base]
\label{q3:static:base-completion}
The positive set \(\mathcal P\) is connected; write
\(\mathcal E=\mathcal P\). It contains the AF end and a deleted collar
of every original boundary component. On \(\mathcal E\),
\(\dd A=0\), \(0<\lambda<1\), and
\begin{equation}
\label{q3:static:static-equations}
\lambda\Ric_h=\Hess_h\lambda,\qquad
\Delta_h\lambda=0,\qquad \Scal_h=0.
\end{equation}
Approach to an interior zero of \(N\) has infinite \(h\)-length.
Attaching all components of \(S\) gives a complete base with compact
boundary and
\begin{equation}
\label{q3:static:base-boundary-data}
h|_{TS}=g|_{TS},\qquad
\lambda=0,\qquad
\partial_{\nu_h}\lambda=\kappa,\qquad
\mathrm{II}_h=0.
\end{equation}
At a component with positive boundary lapse, the base metric has a
smooth reflection with odd lapse. At a zero-lapse component the
reflection is \(C^{2,\alpha}\) for every \(\alpha<1\).
The two-copy double is connected, orientable and complete, and its
reflected metric and odd lapse satisfy the scalar and harmonic
equations across every join.
\end{lemma}

\begin{proof}
The hypotheses of Proposition~\ref{prop:static:common-base} were
verified above. It gives the connected positive set, vanishing twist,
static equations, lapse range, boundary attachments, and completeness
with the entire boundary included. Its null-set and horizon cutoff
arguments apply to all the finitely many components; they impose no
common choice of boundary lapse alternative. In particular, the
proposition does not discard any possible complete end approaching
an interior zero of \(N\). Closedness of \(A\) supplies local static
time coordinates, but no global primitive is asserted at this stage.

Take two copies of the attached base and identify corresponding
points of every \(S_i\) by its collar reflection. The result is
connected because the base is connected and its boundary is nonempty.
Giving the second copy the opposite orientation makes the double
orientable. The reflected second jets agree, so the scalar and
harmonic equations hold continuously across each join.

For completeness, fold the double onto the attached base. The fold
is distance nonincreasing, so a Cauchy sequence projects to a Cauchy
sequence in that complete base. A complete locally compact length
space is proper; hence the projected sequence is contained in a
compact base set. Its inverse image in the double is the quotient
of two compact copies and is compact. The original sequence has a
convergent subsequence and therefore converges. This proves
completeness without a classification or simple-connectivity assumption.
\end{proof}

For the later recovery of the original hypersurface, we record the
original-coordinate forms of the two attachments in
Proposition~\ref{prop:static:common-base}. On a component with
\(u|_S>0\),
\begin{equation}
\label{q3:static:simple-norm}
\dd N|_S=2\kappa X^\flat|_S=2u\kappa\,\nu^\flat .
\end{equation}
Use \((N,y^1,y^2)\) as original smooth collar coordinates and write
\[
X^\flat=a(N,y)\dd N+N\beta_A(N,y)\dd y^A,\qquad
 a(0,y)=\frac1{2\kappa}.
\]
Putting \(N=\lambda^2\) gives
\begin{equation}
\label{q3:static:even-base}
\begin{aligned}
h_{\lambda\lambda}&=4\lambda^2g_{NN}+4a^2,\\
h_{\lambda A}&=2\lambda(g_{NA}+a\beta_A),\\
h_{AB}&=g_{AB}+\lambda^2\beta_A\beta_B.
\end{aligned}
\end{equation}
The coefficients on the right are evaluated at \((\lambda^2,y)\);
the diagonal blocks are even and the mixed block is odd. The base
attachment is therefore a homeomorphism to the original closed
collar, though the change \(N=\lambda^2\) is not a boundary
diffeomorphism.

On a component with \(u|_S=0\), let \(s\geq0\) be the original
\(g\)-normal distance. The attachment uses smooth fields \(a,Y\) with
\begin{equation}
\label{q3:static:double-norm}
u=sa,\qquad X=s^2Y,\qquad a|_S=\kappa,\qquad
N=s^2\bigl(a^2-s^2|Y|_g^2\bigr).
\end{equation}
Thus, with \(D=a^2-s^2|Y|_g^2\),
\[
h=g+s^2D^{-1}Y^\flat\otimes Y^\flat,\qquad
\lambda=s\sqrt D,\qquad D|_S=\kappa^2.
\]
Here both fields are smooth in the original one-sided collar.
These formulas also show that the attached base and original exterior
have the same boundary points and induced boundary metric in both cases.

\subsection{Static asymptotics and the compactified end}
\label{q3:static:asymptotic-subsection}

After a constant linear change of the original end coordinates,
\(h=\delta+O_2(r^{-q_0})\) and
\(\lambda=1+O_2(r^{-q_0})\).
Indeed its original constant metric is
\(I+b\otimes b\), by Equation~\eqref{q3:static:asymptotic-field}.
We prove the improvement required for conformal doubling explicitly;
the harmonic-coordinate mechanism agrees with
Bartnik~\cite[Theorem 3.1, Proposition 3.3, and Theorem 4.3]{Bartnik1986}.

\begin{lemma}[Harmonic-coordinate expansion]
\label{q3:static:asymptotics}
There are asymptotically Cartesian coordinates on the end and a
number \(0<\epsilon<1\) such that, with
\(\gamma=\lambda^2h\) and \(U=\log\lambda\),
\begin{equation}
\label{q3:static:gamma-expansion}
\gamma_{ij}=\delta_{ij}+O_k(r^{-1-\epsilon}),\qquad
U=-\frac Mr+\frac{b_1\cdot x}{r^3}
       +O_k(r^{-2-\epsilon})
\end{equation}
for every fixed finite \(k\).
Here \(M>0\); the vector \(b_1\) is unrelated to the shift asymptote
\(b\).
Only the original two metric derivatives and the stated adjoint
decay are needed to start this improvement.
\end{lemma}

\begin{proof}
The three-dimensional conformal formulas applied to
Equation~\eqref{q3:static:static-equations} give
\begin{equation}
\label{q3:static:gamma-equations}
\Ric_\gamma=2\dd U\otimes\dd U,\qquad \Delta_\gamma U=0.
\end{equation}
Initially \(\gamma-\delta,U=O_2(r^{-q_0})\).
We give the coordinate construction with this derivative count.

Move to a sufficiently distant end and extend \(\gamma\) to a
metric on \(\R^3\), equal to \(\delta\) inside a large ball,
using a fixed-ratio cutoff annulus.
In the resulting coordinates write
\[
\Delta_\gamma=(\delta^{ab}+a^{ab})\partial_{ab}+c^a\partial_a .
\]
The coefficients satisfy \(a=O_2(r^{-q_0})\),
\(c=O_1(r^{-1-q_0})\), and their scaled Hölder norms through
the orders needed here are small after increasing the cutoff radius.
The bounded second derivatives in the hypothesis give scaled
Lipschitz bounds on first derivatives, so they suffice for any fixed
Hölder exponent less than one at this stage.

Put \(s_0=1-q_0\in(0,1/2)\).
For a source with weighted scaled \(C^{0,\alpha}\) size
\(O(r^{s_0-2})\), the Euclidean inverse normalized to vanish at zero is
\begin{equation}
\label{q3:static:subtracted-newton}
Pf(x)=-\frac1{4\pi}\int_{\R^3}
\left(\frac1{|x-y|}-\frac1{|y|}\right)f(y)\dd y.
\end{equation}
It maps that weighted norm boundedly to scaled
\(C^{2,\alpha}\) size \(O(r^{s_0})\).
To verify the zeroth-order bound, split at \(|y|=2|x|\).
On the far part, the kernel difference is bounded by
\(C|x||y|^{-2}\), and its integral is \(O(|x|^{s_0})\)
because \(s_0<1\).
On the near part the separate kernels have the same bound because
\(s_0>0\).
Interior Poisson estimates on rescaled annuli give the derivative
and Hölder bounds.
The ordinary local estimates also apply on the fixed inner ball.

The equations
\[
v^i=P\bigl(-c^i-a^{ab}\partial_{ab}v^i-c^a\partial_av^i\bigr)
\]
are contractions in that weighted space.
Indeed the operator norm of the last two terms is bounded by a
constant times the small scaled norms of \(a\) and \(rc\).
They produce harmonic coordinates \(x^i+v^i\) with
\(v=O_{C^{2,\alpha}}(r^{1-q_0})\) and
\(Dv=O(r^{-q_0})\).
These functions are coordinates sufficiently far out.

There is a derivative issue in changing metric coordinates which
must be addressed before using a classical Ricci equation.
The original assumptions bound \(a,c\), respectively, in scaled
\(W^{2,\infty}\) and \(W^{1,\infty}\).
Differentiate the equation for \(v\) once and apply local
\(W^{2,p}\) estimates to \(Dv\) on rescaled annuli.
All differentiated coefficients and forcing are controlled by
these bounds and the preceding \(C^{2,\alpha}\) estimate.
For every finite \(p\) this gives
\[
v=O_{W^{3,p}}(r^{1-q_0}).
\]
Consequently the transformed \(\gamma-\delta,U\) have scaled
\(W^{2,p}\) size \(O(r^{-q_0})\).
Take \(p>3\); Morrey embedding gives scaled \(C^{1,\alpha}\)
control for some \(\alpha>0\).
In the harmonic chart Equation~\eqref{q3:static:gamma-equations}
has the elliptic form
\[
-\tfrac12\gamma^{ab}\partial_{ab}\gamma_{ij}
 +Q_{ij}(\gamma^{-1},D\gamma)=2U_iU_j,\qquad
\gamma^{ab}\partial_{ab}U=0,
\]
where \(Q_{ij}\) is quadratic in \(D\gamma\).
Schauder estimates first give scaled \(C^{2,\alpha}\) control;
differentiating this elliptic system then gives
\(\gamma-\delta,U=O_k(r^{-q_0})\) for each finite \(k\).
No higher original-coordinate derivative bound has been assumed.

We use two elementary exterior Poisson estimates.
If \(w=o(1)\), \(\Delta_\delta w=f\), and
\(f=O_k(r^{-2-t})\), \(1<t<2\), then
\begin{equation}
\label{q3:static:poisson-monopole}
w=\frac ar+O_k(r^{-t}).
\end{equation}
To prove this, cut \(w\) off on a fixed inner annulus and extend it
by zero; this changes \(f\) only on a compact set.
The extended solution equals the Newton potential of its Laplacian,
since their difference is an entire decaying harmonic function.
The coefficient is \(a=-(4\pi)^{-1}\int f\).
For \(|y|<r/2\), subtracting \(1/r\) from the kernel costs
\[
Cr^{-2}\int_{|y|<r/2}|y||f(y)|\dd y=O(r^{-t}).
\]
The remaining region, including the integrable kernel singularity,
has the same bound. Rescaled interior estimates prove the derivative
version.
If instead \(f=O_k(r^{-4-\epsilon})\), \(0<\epsilon<1\), its first
moments converge and the same reasoning gives
\begin{equation}
\label{q3:static:poisson-dipole}
w=\frac ar+\frac{d\cdot x}{r^3}+O_k(r^{-2-\epsilon}).
\end{equation}
Here subtract \(1/r+(x\cdot y)/r^3\) in the near region.
The remainder costs
\(Cr^{-3}\int_{|y|<r/2}|y|^2|f(y)|\dd y
 =O(r^{-2-\epsilon})\); the far region has that order as well.

Set \(\epsilon=2q_0-1\in(0,1)\).
The harmonic-coordinate equations and initial symbol bounds yield
\[
\Delta_\delta(\gamma_{ij}-\delta_{ij}),
\ \Delta_\delta U=O_k(r^{-2-2q_0}).
\]
Equation~\eqref{q3:static:poisson-monopole} gives
\[
\gamma_{ij}=\delta_{ij}+\frac{A_{ij}}r
 +O_k(r^{-1-\epsilon}),\qquad
U=\frac ar+O_k(r^{-1-\epsilon}).
\]
The harmonic-coordinate condition is
\(\partial_j(\sqrt{\det\gamma}\,\gamma^{ij})=0\).
Expanding it at the displayed order gives
\[
\left(A_{ij}-\tfrac12(\tr A)\delta_{ij}\right)
\frac{x^j}{r^3}=O(r^{-2-\epsilon}).
\]
Taking a radial limit shows \(A=\tfrac12(\tr A)I\).
Its trace in dimension three is zero, hence \(A=0\).
Since \(D^2U=O_k(r^{-3})\), the scalar equation now improves to
\[
\Delta_\delta U
=(\delta^{ab}-\gamma^{ab})\partial_{ab}U
=O_k(r^{-4-\epsilon}).
\]
Equation~\eqref{q3:static:poisson-dipole} proves
Equation~\eqref{q3:static:gamma-expansion}, initially with a real \(M\).

Its positivity uses the lapse range from
Lemma~\ref{q3:static:base-completion}.
The function \(-U>0\) is \(\gamma\)-harmonic.
For \(0<\eta<1\) and a fixed positive \(C\), the function
\[
w=r^{-1}+Cr^{-1-\eta}
\]
is positive and strictly \(\gamma\)-subharmonic sufficiently far out:
\(\Delta_\gamma r^{-1}=O(r^{-4-\epsilon})\), while the leading
Laplacian of \(r^{-1-\eta}\) is
\(\eta(1+\eta)r^{-3-\eta}\).
Choose \(c>0\) so small that \(-U\geq cw\) on a large inner sphere.
The difference tends to zero and is superharmonic, so the exterior
minimum principle gives \(-U\geq cw\) throughout that end.
Taking the radial limit proves \(M\geq c>0\).
\end{proof}

The next section completes the finite-component classification and
recovers the original hypersurface. The independent two-sphere argument
in Section~\ref{q3:local-spinor:section} uses only the static-base and
end estimates established here.
